% ----------------------------------------------------------------------
% Parte I, capítulo 2
% ----------------------------------------------------------------------
\chapter{Metadados e representação descritiva}
\label{cap:metadados}

\instrucao{Abra cada capítulo teórico com um parágrafo que diga qual pergunta ele responde. Revise a literatura de forma crítica: compare autores, aponte divergências e diga qual posição a tese adota.}

Este capítulo discute o conceito de metadados e os padrões usados na representação descritiva de documentos em ambiente digital.

\section{Conceito de metadados}
\instrucao{Parta de definições consagradas e chegue à definição operacional usada na tese. Exemplo: metadados como dados estruturados que descrevem um recurso para fins de identificação, acesso e gestão.}

Na tese, metadados são entendidos como dados estruturados que descrevem um recurso para permitir sua identificação, recuperação e gestão. A definição se apoia na tradição da representação descritiva e nos princípios de organização propostos por \citeonline{ranganathan1937}.

\section{Padrões de metadados}
\instrucao{Apresente os padrões em tabela quando a comparação for o objetivo. Indique a fonte de cada informação.}

A Tabela~\ref{tab:padroes} compara os padrões mais usados em repositórios.

\begin{table}[htb]
  \centering
  \caption{Padrões de metadados em repositórios institucionais}
  \label{tab:padroes}
  \small
  \begin{tabularx}{\textwidth}{lXc}
    \toprule
    \textsc{Padrão} & \textsc{Uso típico} & \textsc{Granularidade} \\
    \midrule
    Dublin Core & Descrição mínima para coleta por agregadores & Baixa \\
    MARC 21 & Catálogos de bibliotecas & Alta \\
    RDA & Regras de descrição baseadas em entidades & Alta \\
    \bottomrule
  \end{tabularx}
  \fonte{Elaborada pelo autor com base em \citeonline{iso15836}.}
\end{table}

\section{Modelos conceituais}
O modelo IFLA LRM organiza o universo bibliográfico em entidades como obra, expressão, manifestação e item \cite{riva2017}. A tese adota esse modelo como referência para avaliar a dimensão semântica.
